\begin{figure}[tb]
\centering
\begin{tikzpicture}[x=0.65cm,y=0.7cm,>=Latex,font=\small]
 \draw[->] (0,0)--(16.8,0) node[right] {scale index};
 \foreach \x/\lab in {2/J,4/2J,8/4J,16/8J}
  {\draw (\x,0.08)--(\x,-0.08) node[below] {$\lab$};}
 \draw[thin] (0.55,3)--(16,3);
 \draw[very thick,blue!55!black] (2,3)--(4,3);
 \draw (0.55,2.87)--(0.55,3.13);
 \draw (16,2.87)--(16,3.13);
 \node[above] at (0.8,3.15) {$i(J)=O(\log J)$};
 \node[above] at (8.7,3.08) {history targeted at $J$};
 \draw[thin,->] (0.82,1)--(16.3,1);
 \draw[very thick,blue!55!black] (4,1)--(8,1);
 \draw (0.82,0.87)--(0.82,1.13);
 \node[above left] at (16,1.08) {history targeted at $2J$ (to $16J$)};
 \draw[dashed] (4,0.2)--(4,3.3);
 \draw[->] (4.3,2.78)--(4.3,1.23);
 \node[right] at (4.5,2.05) {compare physical observables};
 \node[below,align=center] at (8,-0.8)
  {Thick segments are retained in the spliced sequence.};
\end{tikzpicture}
\caption{Finite histories overlap well beyond the intervals retained in
the spliced sequence. Coordinates are compared at the same absolute
scale using observables of the original height law. The picture is
schematic near the entry indices; no convergence rate at entry is
assumed.}
\label{fig:histories}
\end{figure}
